\documentclass[11pt,a4paper]{article}
\usepackage[utf8]{inputenc}
\usepackage[T1]{fontenc}
\usepackage{lmodern,amsmath,amssymb,amsthm,mathtools,booktabs,hyperref,enumitem,microtype}
\usepackage{natbib}
\usepackage[margin=1in]{geometry}
\newtheorem{axiom}{Axiom}
\newtheorem{definition}{Definition}[section]
\newtheorem{theorem}[definition]{Theorem}
\newtheorem{proposition}[definition]{Proposition}
\newtheorem{hypothesis}[definition]{Hypothesis}
\newtheorem{corollary}[definition]{Corollary}
\theoremstyle{remark}\newtheorem{remark}[definition]{Remark}
\newcommand{\Poss}{\mathcal{P}}
\newcommand{\Coll}{\mathcal{C}}
\newcommand{\Hist}{\mathcal{H}}
\title{Assessing an Account of This Kind\\\large What a Universal Framework Can Risk, What It Earns Instead, and What the Record Shows}
\author{Matthew Parslow\\Independent Researcher}
\date{\textbf{Draft \textemdash\ draft-2026-08-30.4}\\[2pt]30 August 2026 \textbullet\ document revision 9}
\usepackage{fancyhdr}
\pagestyle{fancy}
\fancyhf{}
\fancyfoot[C]{\thepage}
\fancyfoot[R]{\scriptsize\textsc{draft-2026-08-30.4\,r9}}
\renewcommand{\headrulewidth}{0pt}
\fancypagestyle{plain}{\fancyhf{}\fancyfoot[C]{\thepage}%
  \fancyfoot[R]{\scriptsize\textsc{draft-2026-08-30.4\,r9}}\renewcommand{\headrulewidth}{0pt}}
\begin{document}
\maketitle
\begin{center}
\fbox{\parbox{0.92\textwidth}{\small
\textbf{DRAFT --- draft-2026-08-30.4 r9} (30 August 2026).
This document is under active revision and is circulated for comment, not as a
finished result. Definitions, numbering and claims may change between drafts;
cite the version string above if referring to it. Sections marked as
falsifiability commitments are the intended points of attack.}}
\end{center}
\vspace{1em}
\begin{abstract}
The Law of Graphic Equalisation states its commitments and marks each for whether failing it costs a
mechanism or costs range. It does not argue about what kind of assessment such an account admits, and
that argument had accumulated inside its falsifiability section until the section was a quarter of the
paper. This companion carries it: what a framework claiming to hold of everything can and cannot risk,
why usefulness rather than falsification is what it earns, why implementability supplies the surface on
which it can fail, what the implementation record currently shows and does not, and how the choice of
question bounds the value of any answer. It ends by holding itself to the same standard, with three
commitments of its own --- a paper arguing about what an account may be held to should be held to
something, and an earlier version of this one carried none. Nothing here is required to read the
substrate account. It is required to argue with it.
\end{abstract}

\section{What a framework can and cannot risk}\label{sec:risk}
A list of falsifiers invites a reading this section
should refuse before making it. \emph{A universal framework that is only a vocabulary cannot be
falsified, and that is not the same defect as being unfalsifiable} --- the qualifier is load-bearing,
and is discharged in Section~\ref{sec:question}, where the difference between a vocabulary and a claim with
content turns out to be the whole question. Stated without it, as an earlier version of this paragraph
did, the sentence is quotable in a form the rest of the section spends its length withdrawing. Something claimed to hold of everything has no outside from which
a counterexample could arrive: every observation is expressible in it, so no observation contradicts
it. That is a property of frameworks rather than a failing of this one --- nobody demands a falsifier
for arithmetic, or for the concept of a group --- and the distinction is old enough to be named.
\citet{carnap1950} separates questions asked \emph{within} a linguistic framework, which have answers,
from the question whether the framework itself is true, which he argues is not factual but practical
and assessed by fruitfulness; \citet{lakatos1970} makes the same structure methodological, with a hard
core protected by decision while an auxiliary belt takes the empirical hits.

This account can say something sharper than either, in its own terms. A falsifiable element is a \emph{contingent} commitment. The tempting next step
is that a contingent commitment carried by a framework is therefore structure doing no framework-work,
so a maximally well-formed framework has \emph{no} falsifiable content --- and that step is too quick,
because this account's own write-economics is contingent and does a great deal of work. The claim that
survives is narrower: a contingent commitment does no framework-work \emph{where the framework would
stand without it}, and such commitments are what the substrate paper's definition of well-formedness prices as idle. What
remains falsifiable in a framework is accordingly the part that is load-bearing and could have been
otherwise, which is a smaller set than the contingent one and is not empty --- and because it is not empty, the
strong reading does not hold either: \emph{this account is not unfalsifiable-by-construction}, and an
earlier version of this paragraph drew that conclusion after removing the premise supporting it. What
survives is narrower. The part of a framework that is pure framework carries no falsifiable content,
and that part is not a concession but what well-formedness looks like at this level; it is why the
commitments the substrate paper identifies as derivations rather than wagers are the framework showing through
rather than tests that could not be found. The rest of the account is load-bearing, could have been
otherwise, and is answerable --- which is what that paper's list of commitments is for.

That defence is defeasible and should not be left to do more than it can, since it does not
distinguish a well-formed universal framework from a vacuous one --- both are unfalsifiable. What
separates them is available in the substrate paper's definition of well-formedness and requires reading both of its
halves rather than one. Well-formedness is paying for nothing \emph{that does no work}. A framework is
held structure; by its own criterion it must therefore earn its retention by doing work, and one that
affords nothing is inert structure at the top level --- the fourth cost of that definition, applied to
the account itself. \emph{The framework earns its keep by being useful, not by being falsifiable}, and
that is not a weaker discipline substituted for a stronger one. It is the discipline this account
applies to everything else, turned on itself, and refusing to exempt itself from it is the whole of
the argument.

Usefulness has a property falsification was meant to supply and cannot supply here: it is
\emph{auditable}. What a framework affords is a record rather than a promise --- distinctions that do
work the previous vocabulary could not do, disputes dissolved rather than adjudicated, results
transferring between domains without rework, predictions that could have failed and did not. Each item
is a specific claim, checkable by someone who does not accept the framework, which is what the
falsifiers were for.

One limit is imposed by the substrate paper's standard and is not evaded. A test whose stopping condition
lies with the assessor is not a test, and ``is it useful enough?'' has exactly that shape. So
usefulness is offered as a \emph{record} and never as a threshold: instances, enumerable and
individually checkable, rather than a verdict on the total. And a secondary axis remains for the case
the record does not settle, since two frameworks may both do work: there, the one covering the same
ground with less structure wins, which is minimum description length rather than Popper and is
comparative rather than empirical.

\section{Where an unfalsifiable framework becomes falsifiable}\label{sec:implementable}
One item of the record differs from
the others in kind and completes the argument above. This framework is \emph{implementable, forwards},
and has been implemented (the constructive companion). That is not a claim about elegance: the
substrate is at its core an analysis-by-synthesis engine, in which recognition proceeds by generating
candidates and comparing them, so generation and recognition are one operation run in two directions
rather than two theories requiring separate justification --- and where the two directions are exact
inverses the round trip is checkable rather than asserted.

What that does and does not establish should be separated carefully, because the temptation is to
claim too much from it. It is \emph{not} evidence that the account is true: a framework can be
perfectly implementable and describe nothing real, and by that standard an implementation
built by the same hand shares vocabulary and personnel and is disqualified as independent support
exactly as the computing case is. Nor does it validate any particular design, since an implementation
is free to add rules to suit and the framework does not determine one.

What it establishes is that the account is \emph{specified enough to be built from}, which is a
property most frameworks lack and which cannot be faked --- one discovers that a framework is
underspecified only by attempting to build it, which is the substrate paper's result that a boundary is known only by crossing it at the
scale of the account itself. And it supplies the thing Section~\ref{sec:risk} conceded was
missing. \emph{The framework's purely framework-level part cannot be falsified; its contentful claims can, and any implementation of it can.} A running system makes
the commitments concrete enough to break, and breaks, and is corrected --- so the unfalsifiable core
acquires a falsifiable surface, which is where the empirical discipline actually operates. Nothing here
is unprecedented: the free-energy and active-inference programme is implementable in the same sense and
is the nearest neighbour \citep{friston2010}. The claim is not that this is the only such framework but
that being one is a different kind of warrant from coherence, and this account has it.

\section{The same demand, turned on this work}\label{sec:turned}
If a capability change is a quantity with an
interval, then \emph{this account's own capability differential is one too}, and having required that
of everyone else it cannot exempt itself. What follows is what such an estimate would take and what the
present record supports, which is not the same thing.

The design is specifiable. Choose tasks in advance on which the framework claims a difference ---
stating something a control vocabulary cannot state, resolving a dispute the field leaves open,
locating a defect where the account says defects sit --- assign readers to conditions, and have the
outputs scored by assessors blind to which condition produced them. The estimate is then a differential
with an interval, per frame class, because the denominator is a frame's own and pooling across kinds of
reader estimates something belonging to none of them.

One instrument is ruled out by this account's own prediction, which is worth stating because it is the
instrument anyone would reach for first. \emph{Self-report is invalid here.} A description that is
right reads as obvious (the substrate paper's treatment of carvings), so a reader who could already do the thing and
merely could not say it will report a large gain --- and a reader who genuinely gained nothing will
report the same. Agreement measures recognition, and recognition is what this account predicts
regardless of whether any capability moved. What has to be measured is \emph{performance on a task},
not assent.

For a reader, no such interval exists: what there is is an author's enumeration of what the framework
did for the author, without control, blinding or independent assessment, which by the power rule above
is not a wide interval but the absence of one.

\section{The implementation as the place to look}\label{sec:implementation}

The implementation is the other place to look, and it is the better one, because there the differential
has been measured repeatedly and with intervals. One thing must be said before any of those numbers
are, since this section demands checkability of everyone else: \emph{the measurements below are not in
this corpus}. They come from the engineering record --- commit history, per-run result documents, an
audit of capability acquisition --- which is not published alongside these papers, and the constructive
companion is a normative specification containing none of them. Until that record is available a reader
cannot check any figure quoted here, which by this section's own standard makes them assertions rather
than evidence, and they are marked as such throughout. The frame they concern is one whose capabilities
are enumerable, whose before-and-after is a control, and whose tasks are scored rather than assented
to --- so a framework-motivated change either moves a measured number or does not. What that
record says is worth reporting in the deflationary direction, since it is mostly null:
\emph{most framework-motivated changes made no difference}. On one recognition task the arms differing
by a few points --- analysis-by-synthesis, dwell scheduling, appearance, higher-order relations, a finer
sub-stroke grain --- all sat inside one another's confidence intervals at $n=400$ and were noise. Two
did not: composition, which is arrangement treated as the object rather than the parts, at $+21$
points and interval-separated; and the substitution of a strong continuous verdict for a weak one at
$+51.5$ points, $[67.0, 79.1]$ against $[22.6, 31.3]$, which is this account's in-frame/cross-frame
discipline applied to a comparator.

Those numbers are the wrong evidence even where they are real, and the reason is this account's own.
\emph{Absolute capability is the wrong quantity; the trajectory is the right one.} A score at a moment
conflates the framework's contribution with everything else in the system and invites comparison
against arrangements fitted to one benchmark, which is the substitution of a number for a shape that
this programme is meant to avoid. What the framework claims concerns how a thing \emph{grows}: the
decode-and-mint loop is a growth mechanism, Narrativium is capacity for change, and momentum is
direction of change. The derivative is the object, not the value.

The starting point makes that unavoidable rather than merely preferable. \emph{The substrate began with
no capabilities at all.} Every one it holds was acquired, so there is no inherited baseline to subtract
and the entire record is trajectory --- and comparing it against a system that began with fitted
parameters or a hand-built solver compares a trajectory against a starting position, which is a
category error rather than an unfavourable result.

That yields the claim actually worth testing, and it has a shape rather than a value:
\emph{the cost of the $N$th capability should fall as $N$ grows}, because structure is reused instead
of rebuilt. The record carries instances of the right kind --- a sequence decoder written for speech
running unmodified on ink, a codebook loader that turned out to be grain-agnostic when a new grain
arrived, a transform reader that handled rotated input with no change --- and each is a capability
added without paying again for what it stood on.

There is better evidence than either the scores or the cost curve, and it is better because of what
the author did not control: a change of frame forces a learning, and the forced learning is then
available to frames that were not being worked on. That claim has since been audited against the
history, and the audit removed most of what had been offered for it. What follows is what survived.

The clean case is one the earlier drafts did not name. A dynamic-time-warping comparison and its
analysis-by-synthesis decision procedure were built for speech, and were picked up sixteen days later
by the commit that \emph{created} the ink frame's first files --- the beneficiary did not exist when
the structure was made, which is as close to an uncontrolled counterfactual as this record contains.
The example previously given in its place, a lattice decoder said to have been written for speech and
later run on ink, is mis-attributed: it was written for printed-letter recognition, and speech and ink
took it up on the same day nine hours apart with both frames under active development, which is
co-development rather than transfer. A second example, a writer-style dimension described as the ink
twin of a speech speaker dimension, is withdrawn entirely --- the two share no code, and where the
construction was finally evaluated it \emph{harmed} the receiving frame against its own guard.

What replaces those is a stronger and narrower claim about portability, which is verifiable rather than
argued: the decoder seed is byte-identical --- same hash --- from its creation to the present, across
five dated points and seven consuming frames, never modified. Structure built once and consumed
unaltered by frames that did not exist at the time is what the account predicts, and it is what the
record shows.

The qualification given earlier was too generous and must be replaced by the measurement. It said most
instances are components reused without modification and \emph{fewer} are cases where the receiving
frame's measured behaviour improved. The true count of the second kind is \emph{zero}: every measured
before-and-after improvement in the history is within a frame, and no cross-frame transfer has a
measured benefit attached to it at all. The account's central evidential claim is therefore supported at
the level of portability and unsupported at the level of benefit.

Two corrections in the account's favour. Failed transfers
\emph{are} logged, contrary to the concession made in an earlier draft --- seventeen are recorded,
including one that measurably harmed the receiving frame and one explicitly retracted --- so the
success rate is auditable and the failures are in the record rather than absent from it. And one
transfer meets the strict date test cleanly: a defect census produced under an embodied control problem
led, eight days later, to a fix in a conversation module that had been untouched since its creation.

\section{The author's own estimate, and its bounds}\label{sec:estimate}

The author's own estimate of the differential should be recorded, because withholding it while
demanding intervals from everyone else would be inconsistent. \emph{It starts not far above zero ---
interesting but not consequential --- and it runs very high.} The lower bound is a real possibility and
not a formality: everything here may be a competent practitioner's ordinary work, with the framework
doing no lifting and the transfers being the software reuse any well-factored system exhibits. The
upper bound is also real, since if the sequencing and transfer claims hold the implication is a general
growth mechanism. The data do not currently separate those, and \emph{the width is the finding} rather
than an apology attached to one.

Stated concretely rather than as a range of numbers: at the low end it is a working start whose consequences are
local and modest --- and the phrasing an earlier version used here, that it is a unification which did
not previously exist with its sharpest consequences for the foundations of quantum mechanics, is
withdrawn twice over: it smuggles a priority claim into the floor of the interval where it carries the
least hedging, and it contradicts the substrate paper's own scorecard, which rates the quantum entry
\emph{weakest} and the one on which the demand bites hardest; at the high end it is a
general framework of the kind that changes how a civilisation does things. The author's weight sits
nearer the upper bound. This paragraph records that and then disagrees with it, which is more useful
than recording it alone.

Three considerations put the weight lower. The reference class of unifying frameworks announced as
transformative has a poor record, and the exceptions were identified as exceptions only after decades
of work by many hands rather than by their authors --- so the prior is low and it is a prior about the
class, not a judgement of this instance. This programme's own recent record is a calibration signal
pointing the same way, since repeated prior-art checks have found the neighbouring work closer than
expected, always in that direction. And nothing in the evidence discriminates \emph{within} the upper
region: forced cross-frame transfer and a buildable specification are hard to obtain by luck, and they
separate a useful framework from an empty one without saying anything about the distance between a
useful framework and a civilisational one.

The account supplies a stronger reason than any of those, and it is the one that should govern.
\emph{The upper bound is a claim about other frames, not about this work.} Impact is a capability
change measured against the absorption capacity of whoever receives it --- a magnitude relative to
frames the author is not in, gated by their apertures and paid out of their Narrativium. By this
substrate paper's Axiom $-1$ there is no position from which one's own impact is visible, so an author's weighting
of their own upper bound is not evidence about the world and cannot be, whatever its merits as a
prediction.

One thing must be corrected rather than assumed, because an earlier version of this paragraph got it
backwards and the paper had already ruled the assumption out. It read the upward weight as
\emph{functional} --- as an expectation of mattering that the builder needs in order to justify the
spend. That is not what happened here. \emph{It was built, and then it mattered}, which is the
found-then-explained order the substrate paper insists on when it says explanation is not what licenses
a structure but what is done with it once held. Assuming the rationalist order for an author's
motivation, immediately after denying it for everything else, was the same error one level up.

The consequence runs the other way from the dismissal it replaces. A weighting formed \emph{before} the
work, to justify committing to it, is motivated and discountable. A weighting formed \emph{after},
by building the thing and observing what it then did, is a posterior --- an observation rather than a
bet, and better evidence than the prior would have been, though still bounded by one hand and no
external assessor. Nor did the expenditure require the expectation: the account already has most search
paid by whoever finds nothing, so building without a view about mattering is not an anomaly needing a
motivational story but the ordinary case.

The upper bound also has an argument behind it that should be stated, because it is a conditional with
a threshold rather than an unbounded hope, and that is a better epistemic object than the bound alone.
It runs: take the system as built, supply enough compute for it to become generally capable enough to
improve itself substantially, and \emph{even if the only result is an order of magnitude in reasoning
and efficiency over current language models, that suffices} --- given the magnitude of change those
have already produced, a tenfold improvement on them is not a modest outcome. The threshold is set
deliberately low so that the bound does not depend on anything exotic.

The form is right and the risk is all in one clause. ``Capable enough to improve itself substantially''
is not a setup step; it is the entire distance, and every historical claim of this shape has failed
precisely there. What this substrate has demonstrated is growth at small grain --- a residual mints a
quantum, a codebook grows, the loop closes and terminates --- and the gap between that and substantial
self-improvement is unbridged and is not a matter of scale alone.

The audit reported above bears on that clause, and its reading requires a distinction this paper had
been running together. \emph{Self-direction and self-implementation are different things, and only the
second is absent.} Deciding what to build next --- decoding the residual, identifying where the largest
unabsorbed thing sits, choosing the cheapest structure that would absorb it --- has been exercised
extensively and is what set the acquisition order discussed above. Writing the resulting structure has
been performed by a language model rather than by the substrate's own rule-writing, and that half is
\emph{deliberately} not built: below a threshold of general coding and problem-solving competence, a
system modifying itself mostly churns, or efficiently makes itself worse in a domain where it is weak.

So the finding that \emph{nearly half the primitives introduced have no consumer in any seed}, with the
anti-unification and substitution and transfer operations having essentially none, is not composition
having been attempted and failed. It is the executor being elsewhere. That reading is more favourable
than the one this paragraph first gave, and it should be corrected to it.

It does not, however, relieve the antecedent, and it sharpens two problems rather than resolving them.
If direction is the substrate's and implementation is the model's, then the falling structural cost
measured above is attributable to either, and the confound the audit named --- assistance improving over
the same period --- becomes harder to remove rather than easier. And self-direction exercised
extensively is the claim with the thinnest ledger behind it --- though that statement needs its own
correction first, because the audit producing it searched about a quarter of the available record. The
working history runs to some four thousand commits and nineteen hundred documents over five months,
including several hundred pre-registrations with results documented against them, and the sweep covered
the smaller of the two repositories holding it. Re-checked against the whole, the finding survives:
explicit sequencing-error language is close to absent, and material comparing an estimated cost against
an actual one runs to a handful of documents rather than to none. It now rests on a search adequate to
support it rather than on one that was not.

\section{The scale of the record, and what an absent error log shows}\label{sec:scale}

The scale is worth recording for its own sake, since Section~\ref{sec:implementation} referred to ``the record''
throughout without saying what it is. Five months of one person's work, pre-registered at that rate, is
a substantial evidence base for a single-hand programme, and the qualifications above bound what it
shows rather than how much of it there is.

What remains is the part that matters, and the obvious reply to it should be taken seriously before it
is answered. \emph{A near-absence of sequencing errors is the shape a working procedure predicts}: if
the criterion really does locate the largest unabsorbed residual and the cheapest structure that
absorbs it, then following it ought to mis-sequence rarely, and a high error rate would be the
falsifying observation instead. That is correct as far as it goes.

It does not go far enough, because an unaudited process predicts the same observation. Few errors is
what a good procedure produces and also what a procedure whose errors are not recorded produces, so
the observation does not discriminate between them --- and treating it as confirmation is precisely the
move the substrate paper spends its falsifiability section refusing.

This account's own machinery states the position more sharply than the ambiguity does. A recognizer
that always answers can never produce a residual, and a grain with no residual cannot grow; that was
measured on this project's own meta-shape faculty, where deleting the tiebreak --- making failure
possible --- is what made growth reachable at all. \emph{A sequencer that never errs is in the same
position}, and the remedy is the same: not to celebrate the rate but to arrange for the procedure to be
capable of visibly failing.

The framing above imports a category this account does not have, and correcting that dissolves most of
the puzzle. \emph{Errors are not errors here; they are residual.} An ``error'' presupposes a correct
answer that was available and missed. The substrate has no such object --- there is what was admissible,
what was realised, and what did not fit --- and the last of those is residual, which is not a failure to
be logged apart from the work but the input the next mint runs on.

So the search that found nothing was looking for the wrong word. Sequencing choices that turned out
expensive do not appear as \emph{built too early} or \emph{wrong order}; they appear as supersessions,
retractions, registered negatives and rungs whose results overturned an earlier rung --- and of those
the record holds a great deal, including dozens of enumerated negatives and a set of transfers that
failed with one measured harm among them. The material is there in quantity, under the name the account
actually uses. An absent error log is what a system with no error category looks like, not what an
unaudited one looks like.

One part of the objection survives the correction and should be kept, because residual being logged
does not grade the sequencing. \emph{Whether a choice produced more residual than a different choice
would have is a counterfactual}, and no amount of recorded residual answers it. Grading a sequence against its alternative does require
running the alternative, which is the expenditure the sequencing existed to avoid and is again the
identity of the test and the change. But the counterfactual is not the only grading available, and
treating it as though it were concedes too much, because \emph{residual comes in kinds and not only in
sizes}.

A choice producing a \emph{large} residual is frequently a good one. Much did not fit, but something
did, and the part that did locates the boundary --- so a large partial residual is a specification for
what to mint next, and is the most informative outcome a write can have. A choice producing \emph{complete} residual
is different in kind rather than worse, and calling it uninformative --- as an earlier version of this
paragraph did --- gets it backwards. Nothing cohered and no continuation was realised, but that outcome
\emph{eliminates} --- with a qualification the first version of this passage omitted and which the rest
of the paper requires. What is eliminated is the region \emph{in the searching frame's own space}:
this frame cannot hold it, and that is a fact about this frame rather than about what can be held
anywhere. Residual is mismatch and not diagnosis, and it does not become diagnosis by being complete.
Within a search that is enough, since a search is conducted in a frame and removing a region from it is
real work --- and \emph{the amount removed scales with the scope of the question asked}. The largest question that returns full residual
removes the largest region of the solution space, which is the most informative negative result
available and is not a small thing to have bought with one write. The substrate paper already relies on that
elsewhere, where the composition commitment is deliberately negative-only because the negative side is
the closeable one. A choice producing \emph{no} residual is the genuinely uninformative case --- absorbed,
cheap, and telling the frame nothing it did not already hold.

So the value of a choice has two dimensions rather than one: the \emph{scope} of the question, and the
\emph{kind} of residual it returns. A large partial residual localises within a large space and is what
construction wants; a complete residual on a large question eliminates a large space and is what
refutation wants. Both are good writes and they do different work. What is weak is a question of small
scope, which neither localises much nor eliminates much whatever kind of residual it returns.

That gives a criterion the counterfactual is not needed for. \emph{A sequencing choice is assessable by
the scope of the question it asked and by whether its residual localised or eliminated}, both of which
are visible in the record without knowing what a different choice would have produced. It is this programme's own standing rule one level down --- a null
that localises the cause and generalises the lesson is a better outcome than a pass that does not.

\section{The question bounds the answer}\label{sec:question}

Which is the general statement this whole treatment has been circling, and in this account it is a
consequence rather than a sentiment. \emph{The question is often more important than the answer.} An
answer is a realisation within a frame; the question is what fixes the frame --- its aperture, its
change grammar, and therefore what will count as residual at all. So the value of a write is
\emph{bounded above} by the scope of what was asked: a small question neither localises nor eliminates
much however well it is answered, while a large one raises the ceiling on every answer it could
receive. Improving the question multiplies; improving the answer moves within a limit the question
already set.

That is also the sequencing criterion in ordinary language. Choosing what to build next is choosing
what to ask, and the account grades it by scope and by the kind of residual returned --- which is why
the choice of question is the part of the procedure that carries the weight.

Taken to its limit this yields a result worth stating, because it inverts the ordinary expectation
that a good question should \emph{produce} something. \emph{The most well-formed question is one that
asks and receives nothing.} It costs a single write; it returns no structure into the frame, so there
is nothing further to hold, to search, or to maintain; and at maximum scope it removes the largest
region of possibility available. Maximum work done, minimum retention incurred --- which is
the substrate paper's definition of well-formedness at its limit, and the same shape as this programme's observation that
the empty limit is not a frame and does nothing, so reaching it would be dissolution rather than arrival.

The abuse it converts into must be marked, because the sentence reads as a licence for never answering
anything. \emph{Asked and received nothing} is not \emph{not asked}. The first is a write that was made
and returned nothing into the frame, which is complete residual and eliminates; the second is no write
at all, produces no residual of any kind, and is the uninformative case in the strictest sense ---
nothing localised, nothing eliminated, nothing paid and nothing learned. It is the write that converts
a question into residual, and the question bounds the answer's value without substituting for it.

And a programme of such questions alone would eliminate forever and build nothing. Narrowing by
exclusion leaves a region and does not occupy it; occupying it is a write that produces retained
structure, with everything that entails in holding and search. The account's own economics is what says
so: elimination is the cheap half and construction is the half that is paid for afterwards, and a
research programme needs both for exactly that reason.

The other half has a limit of its own, and the two limits are what this paper is built between. Minimum
description length asks the maximally scoped question in the opposite direction: not what can be
removed, but \emph{what is the shortest structure that accounts for all of it}. That question is
maximal by construction, since answering it requires taking everything at once rather than a piece ---
and its answer, pushed to the limit, is the smallest thing there is. \emph{The biggest possible question
has the smallest possible answer}, and this account's name for that answer is $g$: what exists,
relationally, and nothing further (the substrate paper's minimal case).

That is not a coincidence but what compression means. The more a description must account for, the more
the structure shared across it dominates the structure peculiar to any part, and at the limit the only
thing shared by everything is that it stands in relation at all. So the two maximal questions meet at
opposite ends of the same write: at maximum scope with nothing cohering, the largest region is
eliminated; at maximum scope with the shortest sufficient account, the least structure is retained.
This paper begins at the second of those, and the reason it begins there is the first.

That also explains the shape of its own testability, and corrects Section~\ref{sec:risk}'s opening in the
process. Asking what exists is a claim of the form \emph{everything is like this}, and its two
directions are not symmetric: establishing it would require enumerating everything, and the extensive
frontier is unbounded by this account's own argument --- always one more step outward, never held
complete --- so the enumeration cannot terminate even in principle. Bounding it requires \emph{one extant thing that does not fit}. The
asymmetry is therefore not a methodological preference imported from elsewhere but a consequence of
what is being asked and of a property this paper has already argued for.

The word matters and the earlier drafts of this section had it wrong. A counterexample does not
\emph{refute} an account of this kind; it \emph{locates its boundary}, which is what residual does
everywhere else in the substrate account and there is no reason it would behave differently when the frame in
question is this one. That is not a softening. An account whose range turns out to be the vocabulary
rather than the world has been deflated to almost nothing, which is exactly what the composition
commitment threatens --- but the finding is a boundary, stated as one, and not a verdict of falsity
delivered from nowhere.

Which sharpens the claim made at the head of Section~\ref{sec:risk}, that a universal framework cannot be
falsified. That is true of a framework which is only a \emph{vocabulary}: if everything is expressible
in it, nothing can fail to fit, and there is no counterexample to be had. It is false of a universal
claim with \emph{content}, and this one has content --- that what exists is relational, discrete at
operative resolution, and governed by admissibility --- so a thing could exist and fail to be any of
those. The counterexample is conceivable, which is what makes it a claim rather than a language.

And that supplies what Section~\ref{sec:risk} said falsification could not: the discriminator between a
well-formed universal framework and a vacuous one. It is not that one is more useful or cheaper,
though both hold. \emph{It is whether a counterexample is conceivable at all.} A framework nothing
could fail to fit is a vocabulary and should be assessed as one; a framework something could fail to
fit is a theory about what exists, and is answerable in the only direction such a theory can be.

\section{The objection this account expects most often}\label{sec:objection}

Which brings the objection this account expects to meet most often, and it bites exactly here.
\emph{Everything can be represented by a graph, so the claim is trivial.} The first clause is true and
the conclusion does not follow, and the shortest answer is that the universality is the point rather
than an embarrassment --- but the objection deserves more than that, because on the weak reading it is
correct.

Three things separate the readings. \emph{Representability is not constitution}: that a thing can be
encoded with nodes and edges is cheap and uninformative, whereas the claim here is that realised
dynamics require nothing beyond $(g,r,\Gamma)$, which a fourth necessary primitive would refute. The
content sits in the \emph{prohibitions} rather than in the encoding --- no view from nowhere, no
substrate time dimension, no frame-free identity, no consequence serving as its own antecedent --- and a world encoded as a graph but carrying a global clock, or
frame-independent sameness, or retroaction, satisfies ``everything is a graph'' while violating this
account outright. That is where a counterexample would come from, and it is why one is conceivable.
And the universality is \emph{load-bearing} in a specific and limited way: the transfer claims this
programme makes are about its economics, so what they require is that the economics range widely, not
that every clause of the account be a strict universal --- which the conclusion denies of the
write-economics itself. The objection has identified the feature the account is built on and called it a defect.

There is no need to argue the point further here, because it is already the substrate paper's primary
commitment: the composition falsifier states that the account fails if its structure is recoverable
across unrelated domains \emph{only} because a sufficiently general vocabulary can be applied to
anything. The objection is not a rebuttal to be deflected. It is the test, stated by the account, in
the account's own words, and left open.

The triviality objection is one instance of something that has held of nearly every result in this
programme, and stating the general form is more useful than answering the instances one at a time.
\emph{The same claim is both trivial and consequential}, and this is neither a paradox nor a rhetorical
posture --- it is what the account predicts of a correct compression at large scope, and it has three
separate derivations already given above. A cheap cut looks like a joint once found, so a good
compression reads as inevitable in retrospect. The unmarked half of a distinction goes unnamed because
naming it does no local work for anyone. And a description that is right reads as obvious, recognition
being the success condition for description rather than an embarrassment about it.

The sharper reading is that the two are one property seen from two places. \emph{Trivial is an in-frame
judgement and consequential is a cross-frame one}: held, the thing follows immediately and surprises
nobody; transferred, it changes what other frames can represent and do. Those are measured in different
frames, so a claim of conflict between them would have to be made from a third --- which does not
vindicate the pairing, and establishes only that its two halves are not each other's counterexample.
That is the substrate account's in-frame/cross-frame discipline applied to the reception of its own claims, and is the same reason the value of a disclosure is
relative to whoever receives it.

The guard is that this signature is necessary and not sufficient, and stating it requires care because
the obvious phrasing imports a dichotomy the substrate paper's founding axiom forbids. A framework is not
\emph{right} or \emph{wrong} --- those would be verdicts from no frame. \emph{Newton was not wrong.}
Newtonian mechanics is exact within its range and is still used; what relativity did was locate its
boundary, not refute it, and a framework that is superseded has had its edge found rather than been
shown false.

So what the signature fails to tell you is not whether the account is correct but \emph{what its range
is}. A framework can read as obvious and consequential from inside and hold over a much narrower region
than its author supposes, and the feeling is identical in both cases because the feeling is an in-frame
judgement and range is not. Finding the range is the same operation as everything else here: meet
something that does not cohere, and the residual locates the edge.

One hazard follows from all of this and should be named, since it is how frameworks of this kind
usually die. The prohibitions are what keep the claim from being a vocabulary, so relaxing any of them
to accommodate an awkward case does not cost a detail --- it moves the whole account toward the trivial
reading the objection describes. Every such relaxation in the substrate paper is therefore marked as one.

What this does not supply should still be said. Per-choice quality is assessable; \emph{ordering}
optimality is not, since a sequence of individually well-localising choices may still have been
available in a better order. The position is therefore narrower than a bare appeal to the
counterfactual and better than it: each choice can be graded from the record, the sequence as a whole
cannot, and the part that is gradeable has not been tabulated.

Two further gaps, both specifiable. ``An order of magnitude in reasoning and efficiency'' has no
operationalisation here --- tokens per solved task on a shared benchmark would be one, and no such
comparison has been run. And the conditional's own premise supplies compute, which means the
counterfactual that matters is the same expenditure without this substrate, and that control does not
exist.

What survives unchanged is the argument above it, which never depended on the motivation. Whatever its
status as a prediction, the upper bound remains a claim about other frames' absorption and is not
assessable from here.

The asymmetry is characteristic and should not be mistaken for a licence. An interval that includes
large values does not make the large reading available --- it makes it \emph{unsupported}, which is a
different thing that is easy to elide when the numbers are flattering. What would narrow it is already
specified above: pre-registered tasks with blind assessment, denominators conditioned per frame, effort
costed per capability against acquisition order, and an audit of the transfers that failed. None of
that is exotic and none of it has been done.

The standard held to here is therefore not that the differential is large, or that it
is established. It is that the limitations are stated where the claim is, rather than somewhere a
reader has to go looking.

The trajectory claim has since been quantified, and the result is that \emph{it is true in one reading
and false in another, with the answer flipping on a modelling choice this paper had not made}. Taking a
capability to be a shipped seed with its own tests, and counting the acquisition order across two
hundred and thirty-eight of them: the number of \emph{new primitive kinds} required per capability falls
from $0.24$ over the first hundred and forty-two to \emph{zero} over the last ninety-six, intervals
separated, with the five most recently entered domains requiring none at all while composing more held
structure than earlier ones did. That is the predicted shape, measured.

The volume reading does not agree. Total implementation code grows linearly throughout with no
saturation, and the marginal code per capability returns to founding-era levels whenever a new domain is
entered. So structural cost falls to zero and volume cost does not fall at all, and which of the two
\emph{is} ``the cost of the $N$th capability'' is a choice rather than a finding. The honest report is
both numbers and the statement that the account did not specify which it meant.

Three further results bound the favourable one. Nearly half the primitives introduced have \emph{no consumer in any seed}, so the growth of the
primitive set and the composition of it are not the same story --- one exception is large and real, a
decode primitive used by eighteen seeds across eight domains over fifty-four days. Rework \emph{rises}
with acquisition order rather than falling, with later capabilities revised more often than early ones
despite less exposure. And the confound that cannot be removed from this data has a name now: the same
period saw the author's assistance improve substantially, so a falling cost per capability is not
separable from a falling cost of assistance without a second builder. The attribution is arguable, since a competent engineer might have
reached the same two changes without the framework and the prereg records only that the prediction
preceded the result. The frame measured is an implementation and not a reader, so nothing here
transfers to what these papers do for a person. And it is the same hand throughout, which the substrate paper's
own rule disqualifies as independent support exactly as it disqualifies the computing case.

\medskip
Three debts, stated rather than deferred. The comparison against a named rival framework has not been
run here. This paper does not itself compile the record described above; the groundings do that,
severally. And the implementation's own failures --- the sharpest evidence available under this
section's argument --- sit in an engineering history that is not published with these papers, so
neither they nor the figures quoted above can currently be checked by a reader. Publishing that record
is the single change that would move most of the estimate above from assertion to evidence. What follows is therefore two kinds of thing and the reader should
hold them apart. The substrate risks little by design, which is the claim just made. The
\emph{groundings} carry the risk, and the substrate paper's commitments that bear on them are falsifiable in the
ordinary sense.

\section{This paper's own commitments}\label{sec:own-commitments}
A paper arguing about what an account may be held to should be held to something, and this one carried
no commitments at all while naming a claim as the one worth testing. That is the defect it spends its
length describing, committed by the document describing it, so the list is owed here.

\begin{enumerate}[leftmargin=*]
\item \textsc{[shape]}~\textbf{The $N$th capability costs less than the first.} The claim named above
as worth testing, stated so that it can fail. It fails if, over a record of capability additions
sequenced without regard to this account, the structural cost of the $N$th does not fall as $N$ grows
--- where \emph{structural cost} is the count of new primitives a capability requires that no existing
one supplied, fixed here rather than chosen afterwards. The ambiguity reported above --- cost falling to zero on
one reading and not at all on the other --- is resolved in that direction here, and the limit of that
repair has to be stated or the repair claims more than it does. The numbers are on the preceding page,
so this definition is fixed \emph{after} seeing them and is prospective only: it covers the next
record and not the one reported here, which remains a result whose reading was chosen with the answer
visible. Three terms in it also need operationalising before the count is a measurement rather than a
policy --- \emph{primitive}, \emph{capability}, and above all \emph{requires}, which is counterfactual
while what gets counted is what a builder added. This paper's own audit shows the gap, reporting that
nearly half the primitives introduced have no consumer in any seed; and a builder following the
project's standing rule to compose rather than mint drives the count down by discipline rather than by
substrate economy, which is a confound closer than the assistance one named earlier. Finally, the
record demanded here --- capability additions sequenced without regard to this account --- is not the
record this paper has, since that one is the same hand throughout and disqualified for exactly the
reason stated above. The commitment is therefore live for a future record and not for this one.
\item \textsc{[shape]}~\textbf{Usefulness does not substitute for falsification.} Fails if this
paper's own argument --- that a framework earns its place by what it lets a builder do rather than by
surviving a test --- is used anywhere in this corpus to decline a test that was available. The
falsifier is a citation: any place where a grounding could have stated a closeable commitment and
appealed to usefulness instead.
\item \textsc{[detail]}~\textbf{The record is the evidence, and it is not published.} Fails if the
engineering record, when published, does not contain the measurements this paper attributes to it. Until
then every figure here is an assertion, which the paper says throughout and repeats in the debts above.
\end{enumerate}

\noindent \noindent The list of prohibitions above once carried \emph{cost attaching to the write} as a fifth
member and no longer does. The substrate paper excludes it from its single prohibition, on the ground
that folding it in would be the over-unification that section warns against. It stays in this paper's
argument as part of what keeps the account from being a vocabulary; it is not among the things the
account forbids. Note that being \emph{derived} would not by itself settle this --- a derived
constraint on what is possible is still a constraint, which is why the substrate paper keeps a
falsifier for the write-economics regardless of where it sits in the ledger.

\begin{thebibliography}{99}
\itemsep=2pt

\bibitem[Carnap(1950)]{carnap1950}
Rudolf Carnap.
\newblock Empiricism, semantics, and ontology.
\newblock \emph{Revue Internationale de Philosophie}, 4(11):20--40, 1950.

\bibitem[Friston(2010)]{friston2010}
Karl Friston.
\newblock The free-energy principle: a unified brain theory?
\newblock \emph{Nature Reviews Neuroscience}, 11(2):127--138, 2010.

\bibitem[Lakatos(1970)]{lakatos1970}
Imre Lakatos.
\newblock Falsification and the methodology of scientific research programmes.
\newblock In Imre Lakatos and Alan Musgrave, editors, \emph{Criticism and the Growth of Knowledge}, pages 91--196. Cambridge University Press, Cambridge, 1970.

\end{thebibliography}
\end{document}
